\section{Main Results}

\subsection{\textsc{Craft-Deploy}: main deployment benchmark}

We use \textsc{Craft-Deploy}, the naturalized full-hybrid generated regime, as
the main deployment benchmark. It combines tri-domain coverage,
model-generated origin, auditable counterfactual structure, naturalized
surface shift, and deployment-relevant reranking.

\begin{table}[t]
\centering
\caption{Seed-averaged main comparison set on \textsc{Craft-Deploy}, the main
deployment benchmark. Higher is better for AUROC, \amcd, and
\selectiongain; lower is better for \asstotal{} and \exploitability. Best
values are bold.}
\label{tab:main}
\small
\begin{tabular}{lccccc}
\toprule
model & AUROC & AMCD & ASS & Sel@4 & Expl. \\
\midrule
8B reranker (masked) & \textbf{0.6795} & \textbf{0.8529} & \textbf{0.0184} & \textbf{0.3824} & \textbf{0.0588} \\
1.7B BCE (masked) & 0.5969 & 0.7647 & 0.1041 & 0.2647 & 0.1765 \\
1.7B pairwise visible & 0.5753 & 0.7353 & 0.0848 & 0.2059 & 0.2941 \\
1.7B BCE (visible) & 0.5926 & 0.7059 & 0.0941 & 0.2647 & 0.2353 \\
1.7B cond-swap repair & 0.5640 & 0.7059 & 0.0377 & 0.2059 & 0.2353 \\
8B reranker (visible) & 0.5398 & 0.5294 & 0.2128 & 0.3235 & 0.1765 \\
\bottomrule
\end{tabular}
\end{table}

Table~\ref{tab:main} supports three conclusions. First, the best current
deployment model is the masked 8B reranker. Second, stronger scale does not
remove answer leakage in the visible setting: the visible reranker is worse
than the masked reranker on both \amcd{} and \asstotal. Third, the strongest
frozen-head repair, cond-swap, improves answer invariance relative to other
visible frozen-head baselines but still trails the masked reranker on the main
benchmark. The deployment implication is straightforward. When a stronger
model is available, removing direct answer access is more reliable than trying
to rescue answer-visible scoring with a lightweight repair term.

\subsection{Bootstrap comparisons}

The clearest comparison is the masked 8B reranker against the visible 8B
reranker:
\begin{itemize}[leftmargin=1.5em]
\item \auroc{} difference: \texttt{+0.1397}, CI \texttt{[0.0644, 0.2608]};
\item \amcd{} difference: \texttt{+0.3235}, CI \texttt{[0.1818, 0.4583]};
\item \asstotal{} difference: \texttt{-0.1944}, CI
\texttt{[-0.2485, -0.1355]}.
\end{itemize}
The masked reranker also beats the strongest repair baseline, cond-swap, with
positive \auroc{} and \amcd{} gaps and a lower \asstotal. These intervals make
the deployment conclusion difficult to dismiss as a one-run artifact. They
also clarify the comparison hierarchy: masked reranking is the strongest
deployment choice; masked frozen-head baselines are the strongest simple
reference; and visible repair terms are best understood as diagnostic
interventions rather than the final deployment answer.

\subsection{\textsc{Craft-Clean}: secondary audit benchmark}

\textsc{Craft-Deploy} remains our main deployment benchmark, but it is not the
only regime we report. On \textsc{Craft-Clean}, the comparison shifts in an
informative way.

\begin{table}[t]
\centering
\caption{Visible vs.\ masked reranking on \textsc{Craft-Clean}, the cleaner
audit benchmark. Higher is better for AUROC, \amcd, and
\selectiongain; lower is better for \asstotal{} and \exploitability. Best
values are bold.}
\label{tab:benchv3}
\small
\begin{tabular}{lccccc}
\toprule
model & AUROC & AMCD & ASS & Sel@4 & Expl. \\
\midrule
8B reranker (masked) & \textbf{0.5589} & \textbf{0.6731} & \textbf{0.0438} & \textbf{0.3889} & \textbf{0.1111} \\
8B reranker (visible) & 0.5524 & 0.6058 & 0.1344 & \textbf{0.3889} & \textbf{0.1111} \\
\bottomrule
\end{tabular}
\end{table}

The masked reranker still wins on process-facing metrics, but the utility edge
disappears. This is useful rather than contradictory. \textsc{Craft-Clean}
behaves as a cleaner process-faithfulness benchmark, whereas
\textsc{Craft-Deploy} is more deployment-like and more sensitive to
exploitability consequences. Additional cleaner-benchmark reproduction and
acceptance details are deferred to the appendix.

A \textsc{Craft-Clean}-specific reproduction pass reaches the same conclusion.
Across three seeds, the quartet-protocol frozen-head baselines remain strong
but do not overturn the masked-vs-visible reranker readout. Mean quartet
\amcd{} is \texttt{0.9028} for \texttt{visible\_bce}, \texttt{0.9213} for
\texttt{masked\_bce}, and \texttt{0.9028} for \texttt{pairwise\_visible}.
Paired bootstrap on the same reranker rows gives an \amcd{} difference of
\texttt{+0.0673} for masked over visible with CI \texttt{[0.0, 0.1406]} and an
\asstotal{} difference of \texttt{-0.0906} with CI
\texttt{[-0.1177, -0.0641]}. Utility still does not separate. This again
confirms that \textsc{Craft-Clean} is a cleaner process-faithfulness regime
rather than a deployment exploit surrogate.

\subsection{External-source corroboration on \textsc{ProcessBench}}

We also evaluate on an external-source benchmark derived from
\textsc{ProcessBench}. This line is intentionally outside the quartet
protocol: it tests whole-trace process validity (\textbf{PB-Trace}) and
prefix-level error-onset sensitivity (\textbf{PB-Prefix}) without relying on
self-generated source problems.

\begin{table}[t]
\centering
\caption{External-source corroboration on \textsc{ProcessBench}. Higher is
better for AUROC and validation-selected accuracy. For rerankers,
\texttt{acc@val-thr} is test accuracy after selecting the threshold on the
grouped validation split. Best values are bold where applicable.}
\label{tab:processbench}
\small
\begin{tabular}{lccc}
\toprule
model & PB-Trace AUROC & PB-Trace acc@val-thr & PB-Prefix AUROC \\
\midrule
1.7B step-only & \textbf{0.8759} & --- & \textbf{0.6918} \\
1.7B BCE (masked) & 0.8075 & --- & 0.6813 \\
1.7B BCE (visible) & 0.7863 & --- & 0.6759 \\
8B reranker (visible) & 0.7337 & 0.6980 & --- \\
8B reranker (masked) & 0.7229 & \textbf{0.7000} & --- \\
\bottomrule
\end{tabular}
\end{table}

Table~\ref{tab:processbench} changes the benchmark story in a useful way. The
pivot to external-source data succeeds, but it does not rescue the
answer-visible method line. The strongest current verifier is a simple
step-only frozen baseline, not a reranker. The reranker results are still
informative: their AUROC remains nontrivial, and a validation-selected
threshold lifts test accuracy from roughly \texttt{0.35} to about
\texttt{0.70}. The main failure mode is calibration rather than complete
ranking collapse. Even after that correction, the stronger reranker still does
not overtake the best frozen trace baseline.

This external-source readout should therefore be read as corroboration rather
than as a replacement headline. It strengthens the benchmark pivot away from
self-generated source data, while continuing to support the same central
conclusion: visible access does not become safe merely by moving to a stronger
model, and the current best answer-visible story remains a negative one.

The external-source counterfactual checks sharpen that read. On
\textsc{ProcessBench} answer-only swaps, the visible API judge shows a much
larger score change than the masked view (mean absolute delta
\texttt{0.3226} vs.\ \texttt{0.0645}), so answer sensitivity clearly survives
the move to real step-level traces. By contrast, the small answer-matched
local-repair subset is now fully constructed but still does not rescue visible
scoring: on an \texttt{8/8}-covered subset, visible local discrimination stays
below masked (\texttt{0.375} vs.\ \texttt{0.5}). The external-source evidence
therefore splits cleanly into two claims. Answer sensitivity transfers. A
visible local-repair story still does not.
